\clearpage\begingroup\SourceLeaf{441b}\ArithmeticHead{110}{L}{STÄCKEL, GAUSS ALS GEOMETER.}
\renewcommand{\thefootnote}{\arabic{footnote})}\setcounter{footnote}{1}
\begin{center}
\textbf{33.}\par\smallskip
\textls[100]{Bedeutung und Wirkung der Disquisitiones generales.}
\end{center}
\noindent\hspace*{\parindent}In den \textit{Disquisitiones generales} wird nur ein Geometer mit Namen erwähnt:
\textsc{Euler}. Fast alles, was dieser über die Krümmung der Oberflächen gelehrt
habe, sagt \textsc{Gauss} im art. 8 der \textit{Disq.}, sei in den von ihm gegebenen Sätzen
I bis IV enthalten; augenscheinlich sind \textsc{Eulers} 1763 verfasste \textit{Recherches
sur la courbure des surfaces}\footnote{Histoire de l’Acad., année 1760, Berlin 1767, Mémoires S. 119.} gemeint. Die Untersuchungen von \textsc{Gauss} berühren
sich aber noch in einer Reihe anderer Punkte mit denen \textsc{Eulers}, und
wenn es auch unentschieden bleiben muss, ob \textsc{Gauss} die betreffenden Abhandlungen
gekannt hat oder nicht, so scheint es doch um so mehr angebracht,
die Berührungspunkte festzustellen, als dadurch die Fortschritte, die wir \textsc{Gauss}
verdanken, in ein helleres Licht treten.
\par
Es möge zunächst an die in den vorangehenden Nummern erwähnten Arbeiten
\textsc{Eulers} zur konformen Abbildung, über die kürzesten Linien und über
die Abwicklung krummer Flächen auf die Ebene erinnert werden. Für die
kürzesten Linien kommen ausser der grundlegenden Abhandlung vom Jahre
1729 noch zwei Veröffentlichungen in Betracht. In der einen vom Jahre
1755 hatte \textsc{Euler} die Anfänge einer sphäroidischen Trigonometrie entwickelt,
einer Lehre von den Dreiecken, deren Seiten kürzeste Linien eines Drehellipsoides
sind; auch hatte er vorgeschlagen, dass man solche Dreiecke in
der Geodäsie benutzen solle\footnote{\textsc{L. Euler}, \textit{Éléments de la trigonométrie sphéroïdique tirés de la méthode des plus grands et plus petits}, Histoire de l’Acad., année 1753, Berlin 1755, Mémoires S. 258.}. In der zweiten, erst 1806 gedruckten Ab-
\clearpage\SourceLeaf{442}\ArithmeticHead{111}{R}{DIE ALLGEMEINE LEHRE VON DEN KRUMMEN FLÄCHEN.}\setcounter{footnote}{0}
\noindent handlung, die am 25. Januar 1779 der Petersburger Akademie vorgelegt worden
war, kommt er auf die allgemeine Lehre von den kürzesten Linien zurück
und stellt deren Differentialgleichungen für den Fall auf, dass die krumme
Fläche durch irgend eine Gleichung zwischen den kartesischen Koordinaten
gegeben wird, während man früher immer vorausgesetzt hatte, dass die Gleichung
nach einer Koordinate aufgelöst sei.
\par
Die Einsicht, dass die drei kartesischen Koordinaten gleichberechtigt sind,
kommt bei \textsc{Euler} aber auch dadurch zum Ausdruck, dass er bei den Untersuchungen
über die Abwicklung krummer Flächen die drei Koordinaten sogleich
als Funktionen zweier Hilfsgrössen ansetzt. Wie \textsc{Kommerell} mit Recht
bemerkt\footnote{\textsc{M. Cantor}, \textit{Vorlesungen über Geschichte der Mathematik}, Bd. IV, Leipzig 1908, Abschnitt XXIV: \textsc{Kommerell}, \textit{Analytische Geometrie des Raumes und der Ebene}, S. 529.}, liegt hierin der erste Schritt zu der Auffassung der krummen
Flächen als selbständiger Gebilde, die erst \textsc{Gauss} mit vollem Bewusstsein
ihrer Bedeutung durchgeführt hat. Ebenso hat \textsc{Gauss}, geleitet von dem allgemeinen
Begriff der Abbildung, jene Parameterdarstellung zur Grundlage seiner
allgemeinen Untersuchungen über die krummen Flächen gemacht.
\par
Endlich ist eine 1775 verfasste, 1786 gedruckte Arbeit über Raumkurven\footnote{\textsc{L. Euler}, \textit{Methodus facilis omnia symptomata linearum curvarum non in eodem plano sitarum investigandi}, Acta Petrop., t. 6 pro anno 1782: I, 1786, S. 19, 37.}
zu erwähnen, in der \textsc{Euler} die Eigenschaften solcher Kurven in der Umgebung
eines Punktes untersucht, indem er durch den Mittelpunkt der Einheitskugel
Parallelen zu den Tangenten zieht, ganz ähnlich wie \textsc{Gauss} im art. 2 der
Neuen allgemeinen Untersuchungen bei ebenen Kurven den Einheitskreis verwendet.
Bei \textsc{Gauss} findet sich, wie schon erwähnt wurde, die Beziehung der
Richtungen im Raume auf die Punkte der Einheitskugel schon in einer auf
das Ende des Jahres 1799 zu setzenden Notiz (Scheda Ac, Varia, begonnen
Nov. 1799, S. 3). In der Selbstanzeige der \textit{Disq. gen.} sagt \textsc{Gauss}: »Dies Verfahren
kommt im Grunde mit demjenigen überein, welches in der Astronomie
in stetem Gebrauch ist, wo man alle Richtungen auf eine fingierte Himmelskugel
von unendlich grossem Halbmesser bezieht« (W. IV, S. 342); man darf
daher annehmen, dass der Gedanke der Abbildung auf die Einheitskugel
(Himmelskugel) der Astronomie seinen Ursprung verdankt.
\par
Die Abbildung einer krummen Fläche auf die Einheitskugel mittels
\clearpage\SourceLeaf{443}\ArithmeticHead{112}{L}{STÄCKEL, GAUSS ALS GEOMETER.}\setcounter{footnote}{0}
\noindent paralleler Normalen ist schon vor \textsc{Gauss} betrachtet und mit der Lehre von
den Doppelintegralen in Zusammenhang gebracht worden, und zwar von \textsc{O. Rodrigues}
in einer 1815 veröffentlichten Abhandlung\footnote{\textsc{O. Rodrigues}, \textit{Sur quelques propriétés des intégrales doubles et des rayons de courbure des surfaces}, Correspondance sur l’école polytechnique, t. 2, 1815, S. 162; abgedruckt im Bulletin de la société philomatique, année 1815, S. 34; vgl. \textsc{P. Stäckel}, \textit{Bemerkungen zur Geschichte der geodätischen Linien}, Leipziger Berichte 1893, S. 456.}, ganz ähnlich wie es \textsc{Gauss}
selbst in einer Notiz über die Oberfläche des dreiachsigen Ellipsoides tut,
die wohl bald nach 1813 verfasst ist (W. VIII, S. 367). \textsc{Rodrigues} hat auch
schon erkannt und genau auf dieselbe Weise wie \textsc{Gauss} im art. 7 der \textit{Disq.
gen.} bewiesen, dass das Verhältnis der Abbildung eines Flächenelementes auf
die Einheitskugel zu dem Flächenelement gleich dem Produkte der zugehörigen
Hauptkrümmungen ist. Er folgert daraus, dass das Doppelintegral, das \textsc{Gauss}
als Gesamtkrümmung eines Flächenstückes bezeichnet hat, den Inhalt der
Area auf der Kugelfläche angibt, die durch jene Abbildung erhalten wird, und
da einer geschlossenen Fläche die ganze, einfach oder mehrfach bedeckte
Oberfläche der Kugel entspricht, so ergibt sich der Wert des zugehörigen
Doppelintegrales gleich einem positiven oder negativen Vielfachen von \(2\pi\).
\par
Auf einen zweiten Geometer wird in den \textit{Neuen allgemeinen Untersuchungen}
und in den \textit{Disq. gen.} hingedeutet. Auf \textsc{Monge} bezieht sich nämlich
die Bemerkung, dass die partielle Differentialgleichung zweiter Ordnung
für die auf die Ebene abwickelbaren Flächen »bisher nicht mit der erforderlichen
Strenge bewiesen war« (W. IV, S. 344; vgl. W. IV, S. 237 und VIII,
S. 437). Dass es sich um \textsc{Monge} handelt, ergibt sich aus dem Briefe an
\textsc{Olbers} vom Juli 1828 (W. VIII, S. 444)\footnote{Der darin erwähnte, »ungezogene Ausfall« von \textsc{Fayolle} steht im Philosophical Magazine, new series, vol. 4, London 1828, S. 436; er ist abgedruckt im Briefwechsel G.-O., 2, S. 508.}; \textsc{Gauss} sagt hier mit Recht, dass
bei \textsc{Monge} das Vorhandensein gerader Linien, nach denen die Fläche gebrochen
wird, erschlichen sei. Im übrigen haben die Untersuchungen des französischen
Geometers, die mehr die Untersuchung besonderer Flächenklassen betreffen,
auf \textsc{Gauss} keinen Einfluss gehabt, und dasselbe gilt auch für dessen Darstellende
Geometrie, die \textsc{Gauss} 1813 mit anerkennenden Worten besprochen
hat (W. IV, S. 359).
\par
Wenn man noch die Anregung hinzunimmt, dass \textsc{Gauss} durch das \textsc{Legendresche}
Theorem über die Zurückführung der kleinen sphärischen Dreiecke
\clearpage\SourceLeaf{444}\ArithmeticHead{113}{R}{DIE ALLGEMEINE LEHRE VON DEN KRUMMEN FLÄCHEN.}\setcounter{footnote}{0}
\noindent auf ebene Dreiecke erfahren hat, so ist alles erschöpft, was sich aus der Zeit
vor 1827 mit seinen Forschungen über die allgemeine Lehre von den krummen
Flächen in Zusammenhang bringen lässt, teils auf deren Gang einwirkend,
teils nur im Strom der Entwicklung auftauchend und wieder untergehend.
\par
Wie gross der Eindruck war, den die \textit{Disq. gen.} sogleich bei ihrem Erscheinen
machten, geht aus den Briefen von \textsc{Bessel} und \textsc{Schumacher} hervor.
Mit den daran anknüpfenden, bedeutenden Arbeiten \textsc{Mindings} beginnt eine
lange Reihe von Arbeiten, deren Ausgangspunkt die Untersuchungen von \textsc{Gauss}
bilden. Es muss jedoch hier genügen, einige noch nicht erwähnte Abhandlungen
herauszugreifen, die in besonders engen Beziehungen zu den \textit{Disq. gen.}
stehen, und die Wirkung der Grundgedanken auf die weitere Entwicklung in
aller Kürze zu schildern; dabei soll die Geodäsie ganz aus dem Spiele bleiben
und für sie auf den schon erwähnten Aufsatz von \textsc{Galle} verwiesen werden.
\par
Schon \textsc{Euler} hatte in seiner Abhandlung über die Krümmung der Flächen
nach einem passenden Masse (juste mesure) für die Krümmung solcher Gebilde
gefragt, einem Masse, das sich der Krümmung der Kurven an die Seite stellen
lasse, und unter Hinweis auf die Sattelflächen erklärt, dass es auf diese Frage
keine einfache Antwort gebe; man müsse vielmehr die Gesamtheit der Krümmungen
in Betracht ziehen, die den zu einem Punkte gehörigen Normalschnitten
zukommen\footnote{Diese richtige Einsicht hat \textsc{Euler} nicht davor bewahrt, bald darauf, 1769, in der \textit{Dioptrica} (Lib. I, § 4, Opera omnia, ser. 3, vol. 3, S. 8) zu behaupten, ein Flächenelement lasse sich stets als sphärisch ansehen, und damit in einen Fehler zurückzufallen, den schon \textsc{Leibniz} begangen hatte (Brief an \textsc{Joh. Bernoulli} vom 29. Juli 1698, \textit{Commercium epistolicum}, Lausanne und Genf 1745, t. 1, S. 387, \textsc{Leibnizens} \textit{Mathematische Schriften}, herausgegeben von \textsc{C. J. Gerhardt}, 1. Abt., Bd. 3, Halle 1855, S. 526). Auch \textsc{D’Alembert} hat sich dieses Fehlers schuldig gemacht (Artikel \textit{Surfaces courbes} in der Encyclopédie méthodique, Abteilung Mathematik, Bd. II, Paris 1784, S. 464).}. Später war bei Untersuchungen über biegsame Flächen,
besonders über die Gestalt von Flüssigkeitshäutchen, das arithmetische Mittel
der beiden Hauptkrümmungen aufgetreten, das schon in der von \textsc{Lagrange}
(1765) begründeten Lehre von den Minimalflächen eine Rolle spielte. \textsc{Sophie
Germain} hat dafür 1831 den Ausdruck mittlere Krümmung vorgeschlagen\footnote{\textsc{S. Germain}, \textit{Mémoire sur la courbure des surfaces}, \textsc{Crelles} Journal, Bd. 7, 1831, S. 1.};
in einem Briefe an \textsc{Gauss} vom 28. März 1829 bemerkt sie, dieser verfahre
geometrisch, sie selbst mechanisch, denn die elastische Kraft, welche die Fläche
\par\smallskip\noindent{\small X\textsubscript{2} Abh. 4.}\hfill{\small 15}\par
\clearpage\SourceLeaf{445}\ArithmeticHead{114}{L}{STÄCKEL, GAUSS ALS GEOMETER.}\setcounter{footnote}{0}
\noindent in ihre ursprüngliche Gestalt zurücktreibt, sei der mittleren Krümmung proportional
(Brief im Gauss-Archiv). Nach \textsc{Sturm}\footnote{\textsc{R. Sturm}, \textit{Ein Analogon zu Gauss’ Satz von der Krümmung der Flächen}, Mathemat. Annalen, Bd. 21, 1883, S. 379.} lässt sich die mittlere
Krümmung auf eine ähnliche Art wie das \textsc{Gauss}sche Krümmungsmass erklären;
beschreibt man nämlich um einen Flächenpunkt eine Kugel und bildet
die in die Fläche eingeschnittene Kurve mittels paralleler Normalen auf die
Einheitskugel ab, so ist der Grenzwert des Verhältnisses der Umfänge beider
Kurven gleich der mittleren Krümmung. Später hat \textsc{Casorati}\footnote{\textsc{F. Casorati}, \textit{Mesure de la courbure des surfaces suivant l’idée commune}, Acta math. 14, 1890, S. 95; vgl. auch \textsc{R. v. Lilienthal}, \textit{Zur Theorie des Krümmungsmasses der Flächen}, ebenda, 16, 1892, S. 143.} das Wort
»Krümmung« beanstandet, weil man auch den Flächen vom \textsc{Gauss}schen Krümmungsmasse
Null eine gewisse Krummheit zuschreiben müsse, und als ein der
Anschauung besser entsprechendes Mass das arithmetische Mittel der Quadrate
der Hauptkrümmungen vorgeschlagen. »Demgegenüber ist zu bemerken, dass
es für eine Fläche überhaupt keinen Ausdruck geben kann, der dem für die
Krümmung einer Kurve völlig entsprechend und zugleich erschöpfend wäre.
Es lassen sich vielmehr von verschiedenen Gesichtspunkten aus für die Flächenkrümmung
mehr oder minder kennzeichnende Ausdrücke aufstellen, die ebenfalls
als Grenzwerte anzusehen sind«\footnote{\textsc{R. v. Lilienthal}, \textit{Die auf einer Fläche gezogenen Kurven}, Encyklopädie der mathematischen Wissenschaften, Bd. III, Teil 3, S. 172 (1902).}. Jedenfalls hat sich unter ihnen der
\textsc{Gauss}sche Ausdruck durch die Fruchtbarkeit seiner Anwendungen ausgezeichnet.
\par
Im Laufe der Zeit hat sich immer klarer die Wichtigkeit der Formeln
im art. 11 der \textit{Disq. gen.} herausgestellt, vermöge deren die zweiten Ableitungen
der kartesischen Koordinaten eines Punktes der Fläche als lineare homogene
Funktionen der ersten Ableitungen und der Richtungscosinus der Normalen
dargestellt werden. \textsc{Weingarten} hat gezeigt, wie man aus ihnen fast unmittelbar
die bei dem Biegungsproblem auftretende partielle Differentialgleichung
zweiter Ordnung für eine der kartesischen Koordinaten ableiten kann\footnote{\textsc{J. Weingarten}, \textit{Über die Theorie der auf einander abwickelbaren Oberflächen}, Festschrift der Technischen Hochschule zu Berlin, 1884.}. Auf
dem von \textsc{Gauss} gebahnten Wege weitergehend, haben \textsc{Mainardi}\footnote{\textsc{Mainardi}, \textit{Su la teoria generale delle superficie}, Giornale dell’Istituto lombardo, t. 9, 1857, S. 394.} und \textsc{Codazzi}\footnote{\textsc{D. Codazzi}, \textit{Sulle coordinate curvilinee d’una superficie e dello spazio}, Ann. di mat. (2), 1, 1867, S. 293; 2, 1868, S. 101, 269; \textit{Mémoire relatif à l’application des surfaces les unes sur les autres}, Mém. prés. par divers sav., 2. série, t. 27, Paris 1883 (vorgelegt 1859).}
\clearpage\SourceLeaf{446}\ArithmeticHead{115}{R}{DIE ALLGEMEINE LEHRE VON DEN KRUMMEN FLÄCHEN.}\setcounter{footnote}{0}
\noindent der \textsc{Gauss}schen Gleichung zwischen den Fundamentalgrössen erster und zweiter
Ordnung zwei Gleichungen hinzugefügt, in denen auch noch die ersten partiellen
Ableitungen der Fundamentalgrössen zweiter Ordnung auftreten, und
\textsc{Bonnet}\footnote{\textsc{O. Bonnet}, \textit{Mémoire sur la théorie des surfaces applicables sur une surface donnée}, Journal de l’école polytechnique, t. 25, cah. 42, 1867, S. 31.} hat bewiesen, dass umgekehrt durch solche Fundamentalgrössen
erster und zweiter Ordnung, die den drei Fundamentalgleichungen genügen,
die Fläche, abgesehen von ihrer Lage im Raume und einer Spiegelung, vollständig
bestimmt wird.
\par
Schliesslich mögen noch Untersuchungen erwähnt werden, die bei Lebzeiten
von \textsc{Gauss} angestellt worden sind und eine Verallgemeinerung seines
Lehrsatzes über die Winkelsumme eines geodätischen Dreiecks bezweckten.
\textsc{Jacobi}\footnote{\textsc{C. G. J. Jacobi}, \textit{Demonstratio et amplificatio nova theorematis Gaussiani de curvatura integra trianguli in data superficie e lineis brevissimis formati}, \textsc{Crelles} Journal, Bd. 16, 1837, S. 344; Werke, Bd. 7, S. 26.} hat im Jahre 1836 den Satz auf Dreiecke ausgedehnt, die von beliebigen
Raumkurven gebildet werden, wobei nur vorausgesetzt werden muss,
dass in den Ecken die beiden sich schneidenden Kurven dieselbe Hauptnormale
haben; die Abbildung auf die Einheitskugel erfolgt mittels der Hauptnormalen
der Kurven, die ja bei den geodätischen Linien mit den Normalen der Fläche
zusammenfallen. Er hat dafür einen von dem \textsc{Gauss}schen Lehrsatze unabhängigen,
einwandfreien Beweis gegeben. Bedenklich war jedoch eine Bemerkung,
die er dem Beweis vorausschickte, dass nämlich die Verallgemeinerung
des \textsc{Gauss}schen Lehrsatzes sich ohne Mühe (sine negotio) ergebe, wenn man
beachte, dass jede Raumkurve als geodätische Linie einer gewissen Fläche
angesehen werden dürfe. Dies stimmt zwar für eine einzelne Raumkurve,
allein es ist, wie \textsc{Clausen}\footnote{\textsc{Th. Clausen}, \textit{Berichtigung eines von Jacobi aufgestellten Theorems}, Astron. Nachrichten, Bd. 20, Nr. 457 vom 29. Sept. 1842.} zeigte, im Allgemeinen bereits unmöglich, eine
krumme Fläche zu bestimmen, die zwei sich in einem Punkte schneidende
und dort dieselbe Hauptnormale besitzende Raumkurven als geodätische Linien
in sich fasst\footnote{Vgl. auch die Briefe von \textsc{Schumacher} an \textsc{Gauss} vom 1. Sept., 9. Nov. und 4. Dez. 1842 und dessen Antwort vom 3. Sept. 1842, Br. G.-Sch. IV, S. 82, 92, 101, 83.}. In seiner Erwiderung\footnote{\textsc{C. G. J. Jacobi}, \textit{Über einige merkwürdige Curventheoreme}, Astron. Nachrichten, Bd. 20, Nr. 463 vom 15. Dez. 1842; Werke, Bd. 7, S. 34.} gibt \textsc{Jacobi}, »um einige unbegründeten
\par\smallskip\noindent\hfill{\small 15\textsuperscript{*}}\par
\clearpage\SourceLeaf{447}\ArithmeticHead{116}{L}{STÄCKEL, GAUSS ALS GEOMETER.}\setcounter{footnote}{0}
\noindent Zweifel über die Richtigkeit des Theorems zu beseitigen«, einen vereinfachten
Beweis und bemerkt nebenbei, aus den Darlegungen von \textsc{Clausen} folge, dass
sein Theorem allgemeiner als das \textsc{Gauss}sche sei, womit er stillschweigend jene
Bemerkung (sine negotio) preisgibt.
\par
Wir wenden uns nunmehr zu den Wirkungen, die die \textit{Disquisitiones
generales circa superficies curvas} vom Jahre 1828 im Lauf des neunzehnten
Jahrhunderts ausgeübt haben. Wenn man die gesamte Entwicklung der
mathematischen Wissenschaften während dieses Zeitraums ins Auge fasst, so
sind es zwei Punkte, in denen die Untersuchungen von \textsc{Gauss} zur Flächentheorie
entscheidend eingegriffen haben. Erstens ist \textsc{Gauss}, während man bis
dahin in der Geometrie nur endliche Gruppen von Transformationen betrachtet
hatte, dazu übergegangen, eine unendliche Gruppe (im Sinne von \textsc{S. Lie}) zu
Grunde zu legen, zweitens hat er die Lehre von den krummen Flächen als
die Geometrie einer zweifach ausgedehnten Mannigfaltigkeit in einer Weise
behandelt, die der allgemeinen Lehre von den mehrfach ausgedehnten Mannigfaltigkeiten
den Weg bahnte.
\par
\textsc{F. Klein}\footnote{\textsc{F. Klein}, \textit{Vergleichende Betrachtungen über neuere geometrische Forschungen}, Programm, Erlangen 1872, Math. Annalen, Bd. 43, 1893, S. 67, Gesammelte Mathem. Abhandlungen I, 1921, S. 460.} hat das allgemeine Problem der geometrischen Forschung mit
den Worten formuliert: »Es ist eine Mannigfaltigkeit und in ihr eine Transformationsgruppe
gegeben. Man entwickle die auf die Gruppe bezügliche
Invariantentheorie.« Nachdem die Gruppe der Bewegungen und Spiegelungen
den Ausgangspunkt der geometrischen Forschung gebildet hatte, war man zu
den Gruppen linearer Transformationen übergegangen, die der projektiven
Geometrie eigentümlich sind, und hatte auch andere endliche Gruppen, wie die
der Transformationen durch reziproke Radien, herangezogen. Ein Ansatz zur
Betrachtung unendlicher Gruppen war allerdings schon in der Geometria Situs
gemacht worden, aber die Fragestellung war hier zu allgemein, als dass man
Anhaltspunkte für weitere Untersuchungen hätte gewinnen können; ergeben
sich doch als Invarianten lediglich ganze Zahlen. Dagegen haben die Transformationen
der binären quadratischen Differentialformen zu einer reichgegliederten
Invariantentheorie geführt. Das \textsc{Gauss}sche Krümmungsmass ist
das erste Glied in der Kette solcher Invarianten. Ihm gesellt sich sogleich,
als Beispiel kovarianter Bildungen, die Seitenkrümmung hinzu. Auch findet
\clearpage\SourceLeaf{448}\ArithmeticHead{117}{R}{DIE ALLGEMEINE LEHRE VON DEN KRUMMEN FLÄCHEN.}\setcounter{footnote}{0}
\noindent sich im art. 21 der \textit{Disq. gen.} bei der Lehre von den geodätischen Linien
schon der Differentialparameter erster Ordnung. Für die Weiterführung nach
der Seite der Flächentheorie ist besonders \textsc{Minding} zu nennen\footnote{\textsc{F. Minding}, \textit{Wie sich entscheiden lässt, ob zwei gegebene krumme Flächen auf einander abwickelbar seien oder nicht}, \textsc{Crelles} Journal, Bd. 19, 1893, S. 370.}. Untersuchungen
aus der theoretischen Physik veranlassten \textsc{Lamé}\footnote{\textsc{G. Lamé}, \textit{Leçons sur les fonctions transcendantes et sur les surfaces isothermes}, Paris 1857.} bei krummlinigen
Koordinaten für Punkte des \textsc{Euklidischen} Raumes die Differentialparameter
erster und zweiter Ordnung aufzustellen, und nachdem im Jahre 1867 \textsc{Riemanns}
Habilitationsvortrag vom 10. Juni 1854: \textit{Über die Hypothesen, welche
der Geometrie zu Grunde liegen}, veröffentlicht worden war, hat \textsc{Beltrami}\footnote{\textsc{E. Beltrami}, \textit{Sulla teorica generale dei parametri differenziali}, Memorie dell’Acc. di Bologna, zweite Reihe, Bd. 8, 1869, S. 551, Opere matematiche II, S. 74.} die
allgemeine Lehre von den Differentialparametern quadratischer Differentialformen
mit beliebig vielen Veränderlichen entwickelt. Gleichzeitig damit sind
die Untersuchungen von \textsc{Christoffel}\footnote{\textsc{E. B. Christoffel}, \textit{Über die Transformation der homogenen Differentialausdrücke zweiten Grades}, Journal f. r. u. a. Math., Bd. 70, 1869, S. 46; Gesammelte mathematische Abhandlungen, Bd. I, S. 352.} und \textsc{Lipschitz}\footnote{\textsc{R. Lipschitz}, \textit{Untersuchungen in Betreff der ganzen homogenen Funktionen von \(n\) Differentialen}, Journal f. r. u. a. Mathematik, Bd. 70, 1869, S. 71, Bd. 71, 1870, S. 274, 288, Bd. 72, 1870, S. 1; \textit{Bemerkungen zu dem Prinzip des kleinsten Zwanges}, ebenda, Bd. 82, 1877, S. 316 (im Anschluss an \textsc{Riemanns} 1876 veröffentlichte Pariser Preisarbeit vom Jahre 1861).} über die Transformation
solcher Differentialformen. Damit wurde der Forschung ein Feld erschlossen,
das noch heute nicht abgeerntet ist.
\par
Mit der Verallgemeinerung auf beliebig viele Veränderliche kommen wir
zu dem Gesichtspunkt der mehrfach ausgedehnten Mannigfaltigkeiten.
\par
Für \textsc{Gauss} hatten die mehrdimensionalen Räume eine metaphysische Bedeutung.
Es handelt sich hier um Spekulationen, die im 18. Jahrhundert
weit verbreitet waren und die auch ins 19. hinüberreichen\footnote{Man vgl. etwa \textsc{F. Zöllner}, \textit{Naturwissenschaft und christliche Offenbarung}, Leipzig 1881, sowie die zahlreichen Veröffentlichungen von \textsc{H. Scheffler} in Braunschweig.}. »\textsc{Gauss}, nach
seiner öfters ausgesprochenen innersten Ansicht, betrachtete die drei Dimensionen
des Raumes als eine spezifische Eigentümlichkeit der menschlichen
Seele; Leute, welche dieses nicht einsehen könnten, bezeichnete er einmal in
seiner humoristischen Laune mit dem Namen Böoter. Wir können uns, sagte
er, etwa in Wesen hineindenken, die sich nur zweier Dimensionen bewusst
\clearpage\SourceLeaf{449}\ArithmeticHead{118}{L}{STÄCKEL, GAUSS ALS GEOMETER.}\setcounter{footnote}{0}
\noindent sind; höher über uns stehende würden vielleicht in ähnlicher Weise auf uns
herabblicken, und er habe, fuhr er scherzend fort, gewisse Probleme hier bei
Seite gelegt, die er in einem höhern Zustande später geometrisch zu behandeln
gedächte« (\textsc{Sartorius}, S. 81). Solche Gedanken reichen wohl bis in die Jugend
zurück, denn in dem Briefe an \textsc{Grassmann} vom 14. Dezember 1844 (W. X\textsubscript{1},
S 436) sagt \textsc{Gauss}, dessen Tendenzen in der Ausdehnungslehre begegneten
teilweise den Wegen, auf denen er selbst nun seit fast einem halben Jahrhundert
gewandelt sei; dabei beruft er sich auf die Selbstanzeige vom Jahre
1831, an deren Schluss von »Mannigfaltigkeiten von mehr als zwei Dimensionen«
gesprochen wird (W. II, S. 178). Auch zeigt der Brief \textsc{Wachters} an \textsc{Gauss}
vom 12. Dezember 1816 (W. X\textsubscript{1}, S. 481), dass bei dessen Besuch im April
1816 von Räumen mit beliebig vielen Abmessungen die Rede gewesen war.
\par
Die Äusserung von \textsc{Gauss}, über die \textsc{Sartorius} berichtet hat, fällt in die
Zeit zwischen 1847 und 1855. Dass \textsc{Gauss} sich gerade in den letzten Jahren
seines Lebens eingehend mit mehrfach ausgedehnten Mannigfaltigkeiten beschäftigt
hat, lässt auch eine Stelle in den Beiträgen zur Theorie der algebraischen
Gleichungen vom Jahre 1849 erkennen: »Im Grunde gehört der
eigentliche Inhalt der ganzen Argumentation [beim Beweise der Wurzelexistenz]
einem höheren, von Räumlichem unabhängigen Gebiete der allgemeinen
abstrakten Grössenlehre an, dessen Gegenstand die nach der Stetigkeit
zusammenhängenden Grössenkombinationen sind, einem Gebiet, welches zur
Zeit noch wenig angebaut ist und in welchem man sich auch nicht bewegen
kann ohne eine von räumlichen Bildern entlehnte Sprache« (W. III, S. 79).
In einer bald darauf, im Wintersemester 1850/51, gehaltenen Vorlesung über
die Methode der kleinsten Quadrate hat \textsc{Gauss} Gelegenheit genommen, seinen
Zuhörern einige Gedanken über solche »Mannigfaltigkeiten von mehreren
Dimensionen«, allerdings unter Beschränkung auf die verallgemeinerte Massbestimmung
des \textsc{Euklidischen} Raumes, mitzuteilen (W. X\textsubscript{1}, S. 473—481)\footnote{Vgl. \textsc{P. Stäckel}, \textit{Eine von Gauss gestellte Aufgabe des Minimums}, Heidelberger Berichte, Jahrgang 1917, 11. Abhandlung.}.
\par
Die von \textsc{Gauss} begehrte Lehre von den nach der Stetigkeit zusammenhängenden
Grössenkombinationen hat bekanntlich \textsc{Riemann} in seinem Habilitationsvortrage
vom 10. Juni 1854 begründet; er hat sich dabei für die Konstruktion
des Begriffes einer \(n\)-fach ausgedehnten Mannigfaltigkeit ausdrücklich
\clearpage\SourceLeaf{450a}\ArithmeticHead{119}{R}{DIE ALLGEMEINE LEHRE VON DEN KRUMMEN FLÄCHEN.}\setcounter{footnote}{0}
\noindent auf die vorher genannten Veröffentlichungen von \textsc{Gauss} (\textit{Selbstanzeige} vom
Jahre 1831, \textit{Beiträge zur Theorie der algebraischen Gleichungen} vom Jahre
1849) berufen und bei der weiteren Untersuchung über die in den Mannigfaltigkeiten
waltenden Massverhältnisse als Grundlage die \textit{Disquisitiones generales
circa superficies curvas} bezeichnet. Durch \textsc{Riemann} haben also die Gedanken,
deren Keime sich in der \textsc{Gauss}schen Abhandlung finden, ihre volle
Entfaltung erfahren. In den folgenden Jahrzehnten hat sich die Bedeutung
dieser Gedanken in immer höherem Masse herausgestellt, nicht allein für die
Mathematik, sondern auch für die analytische Mechanik und schliesslich für
die Grundlagen der theoretischen Physik.
\clearpage\endgroup
